\subsection{幺代数}

K 上的幺代数是指一个序偶 \((A,e)\)，其中 A 是一个带单位元的 K 上的代数，而 e 是 A 的单位元。由于 e 由 A 唯一确定，我们有时会滥用语言，径直说 A 是一个幺代数。若 \((A,e)\) 与 \((A',e')\) 是幺代数，则 \((A,e)\) 到 \((A',e')\) 的一个保单位元态射是指一个态射 \(\varphi\)，它把 A 映到 \(A'\)，且使得 \(\varphi(e)=e'\)。\((A,e)\) 的一个幺子代数是指一个序偶 \((A',e)\)，其中 \(A'\) 是 A 的包含 e 的子代数。

我们常将单位元记作 1。

引理 1. --- 设 A 是一个代数。对 A 的每个幂等元 j，A 的子空间 jAj 是使得 xj = jx = x 的 \(x \in A\) 所成的集合。它是 A 的一个子代数，以 j 为单位元。

证明是初等的。

